\section{Problem and Objectives}\label{sec:problem}

\subsection{Problem and Motivation}

An introductory course in Automated Verification during the previous semester
introduced me to temporal logic and its practical role. Lean offered a way to
go beyond running a model-checking algorithm: I could formalise the checker and
also prove that its answer means what I claim. KILT became a substantial test
of that idea across graphs, automata and executable algorithms.

The problem is concrete. Given a finite Kripke structure and an LTL formula,
KILT computes whether the formula holds on every initial infinite execution.
Lean proves that a positive Boolean result is equivalent to this trace
semantics. For a negative result, a separate theorem connects the returned
finite lasso to an actual violating execution. These are kernel-checked claims,
not conclusions drawn only from running examples.

A finite system can still have an infinite bad execution. In the running
example, $\mathbf{G}(\mathit{request}\rightarrow\mathbf{F}\,\mathit{grant})$
requires every request to be granted eventually. A system may enter a waiting
state and repeat its self-loop forever: granting remains reachable but is never
chosen. The finite prefix and repeating loop describe this execution as a lasso.

I did not start from an empty field. My preliminary literature and API audit
found substantial Lean work on LTL and automata, but its scoped search did not
find the executable pieces assembled into the path I needed. I therefore
chose a narrower CSLib-oriented contribution: connect formula compilation,
product emptiness and semantic correctness in one toolkit.

\subsection{Objectives and Deliverables}

I retained the proposal's core objectives and added counterexample extraction
after the verified checker was working:

\begin{enumerate}
  \setlength{\itemsep}{0pt}
  \setlength{\parskip}{0pt}
  \item define LTL formulas and trace semantics, audit relevant Lean libraries,
    and reuse CSLib's automata interfaces where appropriate;
  \item translate formulas to generalized or ordinary B\"uchi automata, then
    combine the property automaton with a finite Kripke structure;
  \item implement executable accepting-run search and prove its correspondence
    with the declared trace semantics;
  \item return a counterexample lasso and prove its projected execution violates
    the supplied formula; and
  \item provide organised Lean source, representative safety and liveness
    examples, and documentation of the design, proofs, trust boundary and
    limitations.
\end{enumerate}

\subsection{Scope, Assumptions and Simplifications}

I restrict the executable checker to finite Kripke structures so that states
and successor sets can be enumerated by the search. The set of initial states is
also finite and nonempty, which ensures that the model has a starting point.

I kept seriality as a separate proof obligation rather than building it into
the model. Partial models and deadlocks therefore remain representable without
silently adding self-loops, while checking an infinite-trace property requires
proof that every state has a successor. Seriality guarantees continuation; it
does not add fairness or force an enabled transition to be selected.

CTL* and frontend automation exceeded this project's scope. I delivered the
verified LTL checker and lasso extraction; both extensions and CSLib
contribution remain future work.
